\begin{titlepage}

    \noindent % Evita la sangría por defecto
    \begin{minipage}[c]{0.45\textwidth}
        % Logo izquierdo (UNC)
        \includegraphics[width=4.5cm]{images/logo_unc.jpg} 
    \end{minipage}%
    \hfill % Este comando empuja los bloques a los extremos
    \begin{minipage}[c]{0.45\textwidth}
        \begin{flushright}
            % Logo derecho (FAMAF)
            \includegraphics[width=6cm]{images/logo_famaf.png} 
        \end{flushright}
    \end{minipage}

\begin{center}

\LARGE
\textbf{Parametrización geométrica y resolución espacial de un microtomógrafo de rayos X de arquitectura abierta}

\vspace{1cm}

por

\vspace{0.3cm}

\textbf{Mateo José Sironi Nuñez}

\vspace{1cm}





\Large
Presentado ante la FACULTAD DE MATEMÁTICA, ASTRONOMÍA, FÍSICA Y COMPUTACIÓN
como parte de los requerimientos para la obtención del grado de 
Licenciado en Física de la

\vspace{0.6cm}

UNIVERSIDAD NACIONAL DE CÓRDOBA

\vspace{0.6cm}

Mayo, 2026

\vspace{1cm}

\textbf{Director}: Dr. Germán Tirao\\

\vspace{1cm}

\normalsize
\doclicenseThis


\end{center}
\end{titlepage}